\subsection{收缩映射与收缩}

定义 4 —— 设 X 是拓扑空间，A 是 X 的一个子集。

我们把 X 到 A 上的收缩映射定义为从 X 到 A 的一个连续映射，它是 A 到 X 中的典范单射的收缩映射。若存在 X 到 A 上的收缩映射，则称 X 可收缩到 A 上，或称 A 是 X 的一个收缩核。

设 X 是拓扑空间，A 是 X 的一个子空间。

设存在 X 到 A 上的一个收缩映射 \(r\)。子空间 A 等同于点 \(x \in \mathrm{X}\)（满足 \(x = r(x)\)）所成的集合。若 X 是分离空间，则 A 在 X 中是闭的。

根据下面的引理，子空间 A 是 X 的收缩核的充要条件是：从 A 到拓扑空间 Y 的每个连续映射都可扩张为从 X 到 Y 的连续映射。

引理 1 —— 设 X 与 Y 是拓扑空间，\(f\colon X\to Y\) 是连续映射。下列条件等价：

\begin{enumerate}
\def\labelenumi{(\roman{enumi})}
\item
  对每个拓扑空间 Z 及每个连续映射 \(g: X \to Z\)，存在连续映射 \(g': Y \to Z\)，使得 \(g = g' \circ f\)；
\item
  映射 f 是单射且有一个连续的收缩映射；
\item
  映射 f 定义一个从 X 到其像 \(f(\mathrm{X})\) 上的同胚，且该像是 Y 的一个收缩核。
\end{enumerate}

设 \(f\) 是单射，\(r: \mathrm{Y} \to \mathrm{X}\) 是映射 \(f\) 的一个连续收缩映射。对每个连续映射 \(g: \mathrm{X} \to \mathrm{Z}\)，映射 \(g' = g \circ r: \mathrm{Y} \to \mathrm{Z}\) 满足 \(g' \circ f = g \circ r \circ f = g\)。这就证明了 (ii)⇒(i)。

设条件 (i) 满足。令 \(g \colon X \to X\) 为恒同映射。由假设，存在连续映射 \(g' \colon Y \to X\)，使得 \(g' \circ f = g\)。由此推出映射 \(f\) 是单射，且 \(g'\) 是 \(f\) 的一个连续收缩映射。

性质 (ii) 与 (iii) 的等价性是显然的。

定义 5 —— 设 X 是拓扑空间，A 是 X 的一个子集。我们把 X 到 A 上的收缩定义为在 A 上固定的一个同伦 \(\sigma\colon X\times I\to X\)，其起点是 X 的恒同映射，其终点是 X 到 A 上的一个收缩映射。若存在 X 到 A 上的收缩，则称 X 可收缩到 A 上，或称 A 是 X 的一个收缩。

注 —— 换言之，X 到 A 上的收缩是具有下列性质的同伦 \(\sigma\colon X \times I \to X\)：

\begin{enumerate}
\def\labelenumi{(\roman{enumi})}
\item
  \(\sigma(x,0)=x\)（对所有 \(x\in X\)）；
\item
  \(\sigma(x,1)\in\mathrm{A}\)（对所有 \(x\in\mathrm{X}\)）；
\item
  \(\sigma(x,t)=x\)（对所有 \(x\in\mathbf{A}\) 及每个 \(t\in\mathbf{I}\)）。
\end{enumerate}

存在 X 到 A 上的收缩的充要条件是：A 到 X 中的典范单射是偶 (A,A) 到偶 (X,A) 上的同伦等价。

设 \(a\) 是 X 的一点。X 到子空间 \(\{a\}\) 上的收缩恰好就是在 \(a\) 处带基点、把映射 \(\mathrm{Id}_{\mathrm{X}}\) 连接到以 \(\{a\}\) 为像的常值映射的同伦。因此，X 可收缩到 \(\{a\}\) 上当且仅当 X 在 \(a\) 处可缩（III, p. 234, 例 3）。

设 X 是拓扑空间，A 是 X 的一个子集。设 \(\sigma\) 是 X 到 A 上的一个收缩。由 \(r(x) = \sigma(x,1)\) 定义的从 X 到 A 的映射 r 是 X 到 A 上的收缩映射，而 \(\sigma\) 是连接 \(\mathrm{Id}_{\mathrm{X}}\) 与 \(r\) 的同伦。于是，关系式 \(r \circ i = \mathrm{Id}_{\mathrm{A}}\) 与 \(i \circ r = r\) 表明映射 \(i\) 与 \(r\) 是同伦等价，且在同伦意义下互逆。

定义 6 —— 沿用前面的记号，称 \(\sigma\) 为强收缩，如果再有 \(r(\sigma(x,t))=r(x)\)（对所有 \(x\in X\) 及所有 \(t\in I\)）。

例 —— 设 X 是 \(B_{n}\) 中原点的余集。把 \(X \times I\) 映到 X 的、由 \((x, t) \mapsto ((1 - t) + t \frac{1}{\|x\|})x\) 给出的映射是 X 到 \(S_{n-1}\) 上的一个强收缩。与之相伴的 X 到 \(S_{n-1}\) 上的收缩映射是由 \(x \mapsto x / \|x\|\) 给出的映射。

引理 2 —— 设 X 是拓扑空间，U 是 X 的开子集，\(\sigma\) 是 X 到 \(X-U\) 上的一个强收缩。那么 \(\sigma(U\times I)\subset\overline{U}\)。特别地，\(\sigma(U\times\{1\})\) 包含在 U 的边界中。

设 \(x \in \mathrm{U}\)。实数 \(t \in \mathbf{I}\)（满足 \(\sigma(x, t) \in \mathrm{U}\)）所成的集合在 \(\mathbf{I}\) 中是开的且包含 0；设 \(s\) 是它的上确界。我们有 \(\sigma(x, s) \in \overline{\mathrm{U}}\)。若 \(s < 1\)，则 \(\sigma(x, s) \notin \mathrm{U}\)（由 \(s\) 的定义）；若 \(s = 1\)，由于 \(\sigma(x, 1) \in \mathrm{X} - \mathrm{U}\)，结论同样成立。于是 \(\sigma(x, s) \in \mathrm{Fr}(\mathrm{U})\)。由强收缩的定义，我们又有 \(\sigma(x, s) = \sigma(\sigma(x, s), 1) = \sigma(x, 1)\)；特别地，\(\sigma(x, 1) \in \overline{\mathrm{U}}\)。因此，\(\sigma(x, 1)\) 属于边界 \(\overline{\mathrm{U}} \cap (\mathrm{X} - \mathrm{U})\)（\(\mathrm{U}\) 的）。

设 \(t \in \mathbf{I}\)；若 \(\sigma(x,t) \notin \mathrm{U}\)，同样有 \(\sigma(x,t) = \sigma(\sigma(x,t),1) = \sigma(x,1)\)，因此 \(\sigma(x,t) \in \overline{\mathrm{U}}\)，引理得证。

